% =====================================================================
% BOZZA — correzioni al focus "risparmio di memoria"
% Da adattare alla tua voce e inserire nei capitoli indicati.
% NON e' testo definitivo: e' una traccia metodologicamente corretta
% per i tre punti della revisione.
% =====================================================================


% ---------------------------------------------------------------------
% PUNTO 1 -> Capitolo "2 - Metodologia" (dichiarazione onesta della misura)
% ---------------------------------------------------------------------
\paragraph{Cosa misura \texttt{memory\_mb}.}
La metrica di memoria riportata e' il \emph{picco di Resident Set Size}
(\texttt{ru\_maxrss}, ottenuto via \texttt{getrusage(RUSAGE\_CHILDREN)}) dell'intero
processo solver, misurato \emph{end-to-end}: copre cioe' grounding, ricerca CDNL
e il propagatore custom in un unico valore. Va dichiarato con precisione che
\textbf{non} si tratta della memoria del solo grounding. Per le varianti lazy
(\texttt{la}, \texttt{lc}, \texttt{la\_aux}, \texttt{la\_co}) l'interprete
SWI-Prolog gira \emph{in-process} nello stesso processo \texttt{clingo}: di
conseguenza il picco RSS misurato include anche l'engine Prolog, il database
asserito a runtime e i file temporanei generati in \texttt{/tmp}. Il confronto
resta corretto e onesto a patto di interpretarlo per quello che e' --- ``picco
di memoria del processo solver'' --- e non come ``memoria del solo grounding''.

Una conseguenza diretta di questa scelta di misura va resa esplicita per non
indurre in errore il lettore: l'engine Prolog ha un \emph{costo fisso} di
inizializzazione, indipendente dalla taglia $n$. Per istanze piccole questo
overhead puo' rendere la RSS delle varianti lazy \emph{superiore} a quella delle
varianti ground-and-solve, nonostante il lazy materializzi ordini di grandezza
in meno di euristiche. Il risparmio di memoria del lazy e' quindi un fenomeno
\emph{asintotico}: emerge quando l'esplosione combinatoria delle euristiche
ground dello standard supera l'overhead costante dell'engine. Nei nostri dati
questo e' gia' visibile alle taglie piccole, dove le curve lazy partono
\emph{sopra} le curve standard nel grafico della RSS.


% ---------------------------------------------------------------------
% PUNTO 2 -> Capitolo "3 - Risultati" (DOVE sta la memoria: la prova)
% ---------------------------------------------------------------------
\paragraph{Dove sta la memoria: il conteggio delle euristiche ground.}
L'argomento centrale della tesi non e' la memoria totale in se', ma \emph{dove}
quella memoria viene spesa. La prova diretta del risparmio e' il numero di
istanze ground delle euristiche in funzione di $n$, confrontando l'approccio
ground-and-solve (\texttt{gc}/\texttt{ga}) con quello lazy (\texttt{lc}/\texttt{la}).

Il meccanismo e' leggibile direttamente dall'encoding. La direttiva standard
\begin{verbatim}
#heuristic b(X) : x(X), not c(X), S = #sum{Y : c(Y)}. [(((n+1)*S)+X), true]
\end{verbatim}
contiene il valore dell'aggregato \texttt{S} \emph{come variabile nel peso}
$((n{+}1)\,S){+}X$. Poiche' gli atomi \texttt{c(Y)} sono atomi di scelta e non
fatti, il grounder non conosce il valore della somma a tempo di grounding:
deve quindi enumerare \emph{tutti} i valori raggiungibili di \texttt{S} e
materializzare una direttiva \texttt{\#heuristic} distinta per ogni coppia
$(X, S)$. Con $X \in \{1,\dots,n\}$ e $S$ che spazia sui valori delle
somme di sottoinsiemi di $\{1,\dots,n\}$ --- cioe' ogni intero da $0$ a
$n(n{+}1)/2$, dato che $1$ appartiene al dominio --- il numero di istanze e'
\[
2 \cdot n \cdot \Big(\tfrac{n(n+1)}{2} + 1\Big) \;=\; \Theta(n^3),
\]
dove il fattore $2$ conta le due regole simmetriche per \texttt{b} e \texttt{c}.
La formula e' verificabile esattamente sui dati: per $n=10$ si ottiene
$2 \cdot 10 \cdot 56 = 1120$, che coincide con il valore misurato nella colonna
\texttt{ground\_heuristics} per \texttt{gc}.

L'approccio lazy elimina alla radice questa enumerazione. L'aggregato
\texttt{aggregate\_all(sum(Y), holds(c(Y)), S)} viene valutato \emph{a tempo di
solving} dal propagatore Prolog, sull'assegnamento parziale corrente, e non a
tempo di grounding su tutti i mondi possibili. Le euristiche materializzate
restano percio' \emph{costanti} ($2$ fatti query-driven, indipendenti da $n$).
Il grafico di \texttt{combined\_heuristics} (e, in forma piu' grezza,
\texttt{ground\_lines}) mostra dunque la curva $\Theta(n^3)$ di \texttt{gc} che
esplode contro la curva piatta di \texttt{la}: \textbf{e' questo conteggio,
non la RSS, a dimostrare il risparmio}. La RSS lo \emph{conferma} sul consumo
reale; il conteggio delle euristiche ground ne \emph{spiega la causa}.


% ---------------------------------------------------------------------
% PUNTO 3 -> Capitolo "4 - Discussione e Conclusione" (caveat onesto)
% ---------------------------------------------------------------------
\paragraph{Il rovescio della medaglia: il costo per-decisione.}
Per onesta' intellettuale va dichiarato che l'approccio lazy non \emph{elimina}
il costo dell'esplosione combinatoria: lo \emph{sposta}. Cio' che lo
ground-and-solve paga una volta sola, in memoria, a tempo di grounding, il lazy
lo paga ripetutamente, in tempo, a ogni decisione del solver, sotto forma di
query all'interprete Prolog in-process. Le metriche
\texttt{total\_prolog\_query\_time\_ms} e, soprattutto, la sua normalizzazione
per-decisione \texttt{prolog\_ms\_per\_decide}
($=\texttt{total\_prolog\_query\_time\_ms}/\texttt{decide\_calls}$) servono
esattamente a quantificare questo trade-off.

La tesi che i dati supportano e' quindi quella \emph{forte ma circoscritta}:
il lazy ottiene un consumo di memoria \emph{sub-lineare} (costante nella parte
euristica) \emph{al prezzo} di un costo in tempo per-decisione che lo
ground-and-solve non sostiene affatto. Non e' invece corretto sostenere che il
lazy ``domina ovunque'': per istanze piccole, dove l'esplosione delle euristiche
e' ancora gestibile, lo ground-and-solve resta competitivo o preferibile, sia
per la RSS (assenza dell'overhead dell'engine Prolog) sia per i tempi (assenza
di query per-decisione). La conclusione corretta individua e discute il punto di
\emph{crossover} oltre il quale il risparmio asintotico del lazy ripaga il suo
costo per-decisione.
